%% ============================================================
\chapter{Minimization --- Solutions}
%% ============================================================

\begin{enumerate}[{3.}1.]

%%------------------------------------------------------------
\item \textbf{Prove $\mu(\DBAR(s,x))=\DBAR_{R_L}(\mu(s),x)$ by induction for the mapping $\mu$ in Theorem 3.1.}

$\mu(s)=[s]_{R^M}$ (the $R^M$-class of state $s$, identified with an $R_L$-class via the homomorphism). Let $P(n)$: $(\forall s\in S)(\forall x\in\Sigma^n)(\mu(\DBAR(s,x))=\DBAR_{R_L}(\mu(s),x))$.

\textit{Base} ($n=0$): $\mu(\DBAR(s,\lambda))=\mu(s)=\DBAR_{R_L}(\mu(s),\lambda)$. \checkmark

\textit{Step}: $x=ya$, $|y|=n$. Using the homomorphism property $\mu(\delta(t,a))=\delta_{R_L}(\mu(t),a)$:
\[
\mu(\DBAR(s,ya))=\mu(\delta(\DBAR(s,y),a))=\delta_{R_L}(\mu(\DBAR(s,y)),a)=\delta_{R_L}(\DBAR_{R_L}(\mu(s),y),a)=\DBAR_{R_L}(\mu(s),ya). \quad\square
\]

%%------------------------------------------------------------
\item \textbf{Show $\DBAR_{E_A}([s]_{E_A},x)=[\DBAR(s,x)]_{E_A}$ by induction.}

\textit{Base} ($x=\lambda$): $\DBAR_{E_A}([s]_{E_A},\lambda)=[s]_{E_A}=[\DBAR(s,\lambda)]_{E_A}$. \checkmark

\textit{Step}: $x=ya$.
\[
\DBAR_{E_A}([s]_{E_A},ya)=\delta_{E_A}(\DBAR_{E_A}([s]_{E_A},y),a)=\delta_{E_A}([\DBAR(s,y)]_{E_A},a)=[\delta(\DBAR(s,y),a)]_{E_A}=[\DBAR(s,ya)]_{E_A}. \quad\square
\]

%%------------------------------------------------------------
\item \textbf{Prove $E_A=\bigcap_{j=0}^\infty E_{jA}$.}

Define $E_{0A}$: $s E_{0A} t\Leftrightarrow(s\in F\Leftrightarrow t\in F)$.

Define $E_{i+1,A}$: $s E_{i+1,A} t\Leftrightarrow s E_{iA} t\wedge(\forall a\in\Sigma)(\delta(s,a) E_{iA} \delta(t,a))$.

$E_A$ (state equivalence): $s E_A t\Leftrightarrow L(A^s)=L(A^t)$ where $A^s$ is $A$ with start state $s$.

($\subseteq$): We show $E_A\subseteq E_{jA}$ for every $j$ by induction on $j$.

\textit{Base} ($j=0$): if $s E_A t$, then $L(A^s)=L(A^t)$. In particular,
\[
\lambda\in L(A^s)\Leftrightarrow \lambda\in L(A^t),
\]
so
\[
s\in F \Leftrightarrow t\in F.
\]
Hence $s E_{0A} t$.

\textit{Step}: assume $E_A\subseteq E_{jA}$ and let $s E_A t$. Then
$L(A^s)=L(A^t)$. For every $a\in\Sigma$,
\[
L(A^{\delta(s,a)})=\{x\in\Sigma^*:ax\in L(A^s)\}
=\{x\in\Sigma^*:ax\in L(A^t)\}
=L(A^{\delta(t,a)}),
\]
so $\delta(s,a) E_A \delta(t,a)$. By the induction hypothesis,
\[
s E_{jA} t
\qquad\text{and}\qquad
\delta(s,a) E_{jA} \delta(t,a)\quad(\forall a\in\Sigma).
\]
Therefore $s E_{j+1,A} t$. Thus $E_A\subseteq E_{jA}$ for every $j$, and so
\[
E_A\subseteq \bigcap_{j=0}^\infty E_{jA}.
\]

($\supseteq$): Suppose $s E_{jA} t$ for every $j$. We claim that for every
$n\geq 0$ and every string $x$ with $|x|\leq n$,
\[
\DBAR(s,x)\in F \Leftrightarrow \DBAR(t,x)\in F.
\]
We prove this by induction on $n$.

\textit{Base} ($n=0$): the only string of length at most $0$ is $\lambda$. Since
$s E_{0A} t$, we have
\[
s\in F \Leftrightarrow t\in F,
\]
that is,
\[
\DBAR(s,\lambda)\in F \Leftrightarrow \DBAR(t,\lambda)\in F.
\]

\textit{Step}: assume the claim holds for $n$. Let $|x|\leq n+1$.
If $x=\lambda$, then $s E_{n+1,A} t$ implies $s E_{nA} t$, so the induction
hypothesis already gives the result. If $x=ay$ with $a\in\Sigma$ and $|y|\leq n$,
then $s E_{n+1,A} t$ implies
\[
\delta(s,a) E_{nA} \delta(t,a).
\]
Applying the induction hypothesis to the pair $\delta(s,a),\delta(t,a)$ and the
string $y$, we get
\[
\DBAR(\delta(s,a),y)\in F \Leftrightarrow \DBAR(\delta(t,a),y)\in F.
\]
Since
\[
\DBAR(\delta(s,a),y)=\DBAR(s,ay)
\qquad\text{and}\qquad
\DBAR(\delta(t,a),y)=\DBAR(t,ay),
\]
the claim follows.

Now let $x\in\Sigma^*$ be arbitrary and take $n=|x|$. Then
\[
\DBAR(s,x)\in F \Leftrightarrow \DBAR(t,x)\in F,
\]
which means
\[
x\in L(A^s)\Leftrightarrow x\in L(A^t).
\]
Therefore $L(A^s)=L(A^t)$, i.e., $s E_A t$. So
\[
\bigcap_{j=0}^\infty E_{jA}\subseteq E_A.
\]
Hence $E_A=\bigcap_{j=0}^\infty E_{jA}$. $\square$

%%------------------------------------------------------------
\item \textbf{Show $\delta_{E_A}$ in Definition 3.8 is well defined.}

$\delta_{E_A}([s]_{E_A},a)=[\delta(s,a)]_{E_A}$. If $[s]_{E_A}=[t]_{E_A}$, i.e., $s E_A t$, then $\delta(s,a) E_A \delta(t,a)$ (proved in Exercise 3.3, $\supseteq$ direction): $L(A^{\delta(s,a)})=\{x:ax\in L(A^s)\}=\{x:ax\in L(A^t)\}=L(A^{\delta(t,a)})$. So $[\delta(s,a)]_{E_A}=[\delta(t,a)]_{E_A}$. $\square$

%%------------------------------------------------------------
\item \textbf{Show $F_{E_A}$ is well defined.}

$F_{E_A}=\{[s]_{E_A}\mid s\in F\}$. If $[s]_{E_A}=[t]_{E_A}$ and $s\in F$: $s E_A t\Rightarrow L(A^s)=L(A^t)\Rightarrow\lambda\in L(A^s)\Leftrightarrow\lambda\in L(A^t)\Rightarrow s\in F\Leftrightarrow t\in F$. So $t\in F$. $\square$

%%------------------------------------------------------------
\item \textbf{Show $\delta^c$ maps into $S^c$.}

$S^c$ = connected states = $\{\DBAR(s_0,x)\mid x\in\Sigma^*\}$. $\delta^c$ is $\delta$ restricted to $S^c$. For any $s\in S^c$ and $a\in\Sigma$: $s=\DBAR(s_0,x)$ for some $x$, so $\delta(s,a)=\delta(\DBAR(s_0,x),a)=\DBAR(s_0,xa)\in S^c$. $\square$

%%------------------------------------------------------------
\item \textbf{Prove Lemma 3.1.}

Suppose $sE_{m+1A}t$. By definition, this means
\[
(\forall x\in\Sigma^* \ni |x|\leq m+1)(\DBAR(s,x)\in F \Leftrightarrow \DBAR(t,x)\in F).
\]
In particular, the same equivalence holds for every $x$ with $|x|\leq m$. Hence
\[
(\forall x\in\Sigma^* \ni |x|\leq m)(\DBAR(s,x)\in F \Leftrightarrow \DBAR(t,x)\in F),
\]
which is exactly $sE_{mA}t$. Therefore $E_{m+1A}\subseteq E_{mA}$. $\square$

%%------------------------------------------------------------
\item \textbf{Prove Lemma 3.3.}

By definition,
\[
sE_{0A}t \Leftrightarrow (\forall x\in\Sigma^* \ni |x|\leq 0)(\DBAR(s,x)\in F \Leftrightarrow \DBAR(t,x)\in F).
\]
The only string of length $0$ is $\lambda$, so
\[
sE_{0A}t \Leftrightarrow \DBAR(s,\lambda)\in F \Leftrightarrow \DBAR(t,\lambda)\in F
\Leftrightarrow s\in F \Leftrightarrow t\in F.
\]
Therefore the equivalence classes of $E_{0A}$ are exactly $F$ and $S-F$, unless one of these sets is empty, in which case there is only the single class $S$. $\square$

%%------------------------------------------------------------
\item \textbf{Prove Theorem 3.9.}

Assume $E_{mA}=E_{m+1A}$. We prove by induction on $k$ that
$E_{m+k,A}=E_{mA}$.

\textit{Base case} ($k=0$): $E_{m+0,A}=E_{mA}$ trivially.

\textit{Inductive step}: assume $E_{m+k,A}=E_{mA}$. Let $s,t\in S$. By Theorem
3.8,
\begin{align*}
sE_{m+k+1,A}t
&\Leftrightarrow sE_{m+k,A}t \wedge
(\forall a\in\Sigma)(\delta(s,a)E_{m+k,A}\delta(t,a)) \\
&\Leftrightarrow sE_{mA}t \wedge
(\forall a\in\Sigma)(\delta(s,a)E_{mA}\delta(t,a))
\end{align*}
by the induction hypothesis. Applying Theorem 3.8 again,
\[
sE_{m+k+1,A}t \Leftrightarrow sE_{m+1,A}t.
\]
Since $E_{m+1A}=E_{mA}$ by hypothesis, this gives
\[
sE_{m+k+1,A}t \Leftrightarrow sE_{mA}t.
\]
Therefore $E_{m+k+1,A}=E_{mA}$. By induction,
$(\forall k\in\NN)(E_{m+k,A}=E_{mA})$. $\square$

%%------------------------------------------------------------
\item \textbf{Prove $\mu(\DBAR_A(s,x))=\DBAR_B(\mu(s),x)$ by induction for a homomorphism $\mu$.}

\textit{Base}: $\mu(\DBAR_A(s,\lambda))=\mu(s)=\DBAR_B(\mu(s),\lambda)$. \checkmark

\textit{Step}: $x=ya$. $\mu(\DBAR_A(s,ya))=\mu(\delta_A(\DBAR_A(s,y),a))=\delta_B(\mu(\DBAR_A(s,y)),a)$ (homomorphism property) $=\delta_B(\DBAR_B(\mu(s),y),a)$ (IH) $=\DBAR_B(\mu(s),ya)$. $\square$

%%------------------------------------------------------------
\item \textbf{Prove $L(A)=L(B)$ for a homomorphism $\mu: A\to B$.}

For any $x\in\Sigma^*$,
\begin{align*}
x\in L(A)
&\Leftrightarrow \DBAR_A(s_{0A},x)\in F_A \\
&\Leftrightarrow \mu(\DBAR_A(s_{0A},x))\in F_B
\qquad\text{(because $\mu(s)\in F_B \Leftrightarrow s\in F_A$)} \\
&\Leftrightarrow \DBAR_B(\mu(s_{0A}),x)\in F_B
\qquad\text{(by Exercise 3.10)} \\
&\Leftrightarrow \DBAR_B(s_{0B},x)\in F_B
\qquad\text{(since $\mu(s_{0A})=s_{0B}$)} \\
&\Leftrightarrow x\in L(B).
\end{align*}
Therefore $L(A)=L(B)$. $\square$

%%------------------------------------------------------------
\item \textbf{Examples with connectivity and $A/_{E_A}$.}

\begin{enumerate}
\item $A$ not connected, $A/_{E_A}$ not connected: take
\[
A=\langle\{\pmb{a}\},\{s_0,s_1,s_2\},s_0,\delta,\{s_0,s_1\}\rangle,
\]
where
\[
\delta(s_0,\pmb{a})=s_1,\qquad \delta(s_1,\pmb{a})=s_1,\qquad
\delta(s_2,\pmb{a})=s_2.
\]
The state $s_2$ is unreachable, so $A$ is not connected. Also $s_0 E_A s_1$,
because from either state every continuation is accepted, while $s_2$ is not
equivalent to them because from $s_2$ every continuation is rejected. Thus
$A/_{E_A}$ has exactly the two states $[s_0]=[s_1]$ and $[s_2]$, and $[s_2]$ is
still unreachable. Hence $A/_{E_A}$ is not connected.

\item $A$ not connected, $A/_{E_A}$ connected: take
\[
A=\langle\{\pmb{a}\},\{t_0,t_1,t_2\},t_0,\delta,\{t_1,t_2\}\rangle,
\]
where
\[
\delta(t_0,\pmb{a})=t_1,\qquad \delta(t_1,\pmb{a})=t_1,\qquad
\delta(t_2,\pmb{a})=t_2.
\]
Again $t_2$ is unreachable, so $A$ is not connected. But now $t_1 E_A t_2$,
because from either state every continuation is accepted. Therefore
$A/_{E_A}$ has just the two states $[t_0]$ and $[t_1]=[t_2]$, and both are
reachable from the start class: $[t_0]$ is the start state and
\[
\delta([t_0],\pmb{a})=[t_1].
\]
Hence $A/_{E_A}$ is connected.
\end{enumerate}

%%------------------------------------------------------------
\item \textbf{Show there exists a homomorphism from $A$ to $A/_{E_A}$.}

Define $\mu: S_A\to S_{E_A}$ by $\mu(s)=[s]_{E_A}$. Then:
\begin{itemize}
\item $\mu(s_{0A})=[s_{0A}]_{E_A}=s_{0,E_A}$. \checkmark
\item $\mu(\delta_A(s,a))=[\delta_A(s,a)]_{E_A}=\delta_{E_A}([s]_{E_A},a)=\delta_{E_A}(\mu(s),a)$. \checkmark
\item $s\in F_A\Leftrightarrow[s]_{E_A}\in F_{E_A}$, i.e., $\mu(s)\in F_{E_A}\Leftrightarrow s\in F_A$. \checkmark
\end{itemize}
So $\mu$ is a homomorphism. $\square$

%%------------------------------------------------------------
\item \textbf{Show there exists a homomorphism from connected $A$ to $A_{R_{L(A)}}$.}

\begin{enumerate}
\item Define $\psi: S\to\{R_{L(A)}\text{-classes}\}$ by $\psi(s)=[x_s]_{R_{L(A)}}$ where $x_s$ is any string with $\DBAR(s_{0A},x_s)=s$ (well-defined by part (b)).

More precisely: since $A$ is connected, for each $s\in S$ there exists $x_s\in\Sigma^*$ with $\DBAR(s_{0A},x_s)=s$. Define $\psi(s)=[x_s]_{R_{L(A)}}$.

\item Well-definedness: if $\DBAR(s_0,x)=\DBAR(s_0,y)=s$, then $x R^A y$ (same state), and $R^A$ refines $R_{L(A)}$, so $x R_{L(A)} y$, giving $[x]_{R_{L(A)}}=[y]_{R_{L(A)}}$. $\square$

\item Homomorphism:
\begin{itemize}
\item $\psi(s_{0A})=[\lambda]_{R_{L(A)}}=s_{0,R_{L(A)}}$. \checkmark
\item $\psi(\delta(s,a))=[x_s a]_{R_{L(A)}}=\delta_{R_{L(A)}}([x_s]_{R_{L(A)}},a)=\delta_{R_{L(A)}}(\psi(s),a)$. \checkmark
\item $s\in F\Leftrightarrow\lambda\in L(A^s)\Leftrightarrow x_s\in L(A)\Leftrightarrow[x_s]_{R_{L(A)}}\in F_{R_{L(A)}}\Leftrightarrow\psi(s)\in F_{R_{L(A)}}$. \checkmark
\end{itemize}
\end{enumerate}

%%------------------------------------------------------------
\item \textbf{Show there may not be a homomorphism $A\to A_{R_{L(A)}}$ if $A$ is not connected.}

Take $\Sigma=\{\pmb{0}\}$ and let
\[
A=\langle \{\pmb{0}\},\{s_0,s_1\},s_0,\delta,\{s_1\}\rangle
\]
where $\delta(s_0,\pmb{0})=s_0$ and $\delta(s_1,\pmb{0})=s_1$. The state $s_1$ is
unreachable, so $A$ is not connected. Since the reachable part is just the
nonfinal loop at $s_0$, we have $L(A)=\emptyset$. Therefore $A_{R_{L(A)}}$ is the
one-state DFA for $\emptyset$, with a single nonfinal state $q$ and
$\delta(q,\pmb{0})=q$.

If there were a homomorphism $\psi:A\to A_{R_{L(A)}}$, then $\psi(s_0)=q$ because
$s_0$ is the start state. But a homomorphism must also preserve finality:
\[
s_1\in F_A \Leftrightarrow \psi(s_1)\in F_{R_{L(A)}}.
\]
This is impossible, because $q$ is the only state of $A_{R_{L(A)}}$ and $q$ is
not final. Hence no such homomorphism exists.

%%------------------------------------------------------------
\item \textbf{Show there may still be a homomorphism $A\to A_{R_{L(A)}}$ if $A$ is not connected.}

Take $\Sigma=\{\pmb{0}\}$ and let
\[
A=\langle \{\pmb{0}\},\{s_0,s_1\},s_0,\delta,\{s_0,s_1\}\rangle
\]
where $\delta(s_0,\pmb{0})=s_0$ and $\delta(s_1,\pmb{0})=s_1$. The state $s_1$ is
unreachable, so $A$ is not connected, but $L(A)=\Sigma^*$ because every string
keeps the start state $s_0$ in a final state.

Therefore $A_{R_{L(A)}}$ is the one-state DFA
\[
\langle \{\pmb{0}\},\{q\},q,\delta',\{q\}\rangle
\]
with $\delta'(q,\pmb{0})=q$. Define $\psi:A\to A_{R_{L(A)}}$ by
\[
\psi(s_0)=q,\qquad \psi(s_1)=q.
\]
This is a homomorphism: the start state is preserved, finality is preserved
because all three states involved are final, and
\[
\psi(\delta(s_i,\pmb{0}))=\psi(s_i)=q=\delta'(q,\pmb{0})
=\delta'(\psi(s_i),\pmb{0})
\]
for $i=0,1$.

%%------------------------------------------------------------
\item \textbf{Show there need not be a homomorphism from $A_{R_L}$ to $A_R$.}

Let $\Sigma=\{\pmb{a}\}$ and let $Q$ be the right congruence with two classes
\[
[\lambda]_Q=\{\lambda\}, \qquad [\pmb{a}]_Q=\{x\in\Sigma^*:|x|\geq 1\}.
\]
Let $R$ refine $Q$ by splitting the nonempty strings according to parity:
\[
[\lambda]_R=\{\lambda\},\qquad
[\pmb{a}]_R=\{x:|x|\text{ is odd}\},\qquad
[\pmb{aa}]_R=\{x:|x|\text{ is even and }|x|\geq 2\}.
\]
Take $L=\{\lambda\}$, which is a union of $Q$-classes. Then
$A_Q$ has two states,
\[
q_0=[\lambda]_Q \quad\text{(start and final)},\qquad q_1=[\pmb{a}]_Q
\]
with $\delta_Q(q_0,\pmb{a})=q_1$ and $\delta_Q(q_1,\pmb{a})=q_1$.

The DFA $A_R$ has three states,
\[
r_0=[\lambda]_R \quad\text{(start and final)},\qquad
r_1=[\pmb{a}]_R,\qquad
r_2=[\pmb{aa}]_R,
\]
with transitions
\[
\delta_R(r_0,\pmb{a})=r_1,\qquad
\delta_R(r_1,\pmb{a})=r_2,\qquad
\delta_R(r_2,\pmb{a})=r_1.
\]

Assume there is a homomorphism $\psi:A_Q\to A_R$. Since start states must map to
start states, $\psi(q_0)=r_0$. Then
\[
\psi(q_1)=\psi(\delta_Q(q_0,\pmb{a}))
        =\delta_R(\psi(q_0),\pmb{a})
        =\delta_R(r_0,\pmb{a})
        =r_1.
\]
But now
\[
\psi(q_1)=\psi(\delta_Q(q_1,\pmb{a}))
        =\delta_R(\psi(q_1),\pmb{a})
        =\delta_R(r_1,\pmb{a})
        =r_2,
\]
which says $r_1=r_2$, a contradiction. Therefore there need not be a
homomorphism from $A_Q$ to $A_R$.

%%------------------------------------------------------------
\item \textbf{Is $\cong$ appropriate to ask if it's a right congruence?}

$\cong$ is a relation on DFAs, not on $\Sigma^*$. Right congruences are relations on $\Sigma^*$. So asking whether $\cong$ (between DFAs) is a right congruence (on strings) is a category error -- it makes no sense. One could ask whether the relation induced on strings by $\cong$ is a right congruence, but the question itself is ill-formed in the original sense.

%%------------------------------------------------------------
\item \textbf{Is $E_A$ a right congruence?}

$E_A$ is a relation on states $S$, not on $\Sigma^*$. So strictly speaking, $E_A$ is not a right congruence (which is a relation on $\Sigma^*$). However, $R^A$ (induced on strings) is a right congruence, and $E_A$ induces $R^A$ via $x R^A y\Leftrightarrow\DBAR(s_0,x) E_A \DBAR(s_0,y)$.

%%------------------------------------------------------------
\item \textbf{If $A$ is reduced then $A^c$ is reduced.}

$A^c$ = connected part of $A$. States of $A^c$ are $S^c\subseteq S$. If $A$ is reduced, no two distinct states in $S$ are equivalent; in particular, no two distinct states in $S^c\subseteq S$ are equivalent. So $A^c$ is reduced. $\square$

%%------------------------------------------------------------
\item \textbf{For a homomorphism $\mu: S_A\to S_B$, prove $\mu(s) E_B \mu(t)\Leftrightarrow s E_A t$.}

($\Rightarrow$) Suppose $\mu(s) E_B \mu(t)$, i.e., $L(B^{\mu(s)})=L(B^{\mu(t)})$. By Exercise 3.10: $L(A^s)=L(B^{\mu(s)})$ and $L(A^t)=L(B^{\mu(t)})$ (since $\mu$ is a homomorphism from $A^s$ to $B^{\mu(s)}$). So $L(A^s)=L(A^t)$, i.e., $s E_A t$.

($\Leftarrow$) Suppose $s E_A t$. Then $L(A^s)=L(A^t)$. By the homomorphism property: $L(B^{\mu(s)})=L(A^s)=L(A^t)=L(B^{\mu(t)})$, so $\mu(s) E_B \mu(t)$. $\square$

%%------------------------------------------------------------
\item \textbf{Define $\psi: M^c\to A_{R^M}$ and prove it is an isomorphism.}

\begin{enumerate}
\item For each state $s\in S^c$, choose a string $x_s\in\Sigma^*$ with
\[
\DBAR(s_0,x_s)=s,
\]
and define
\[
\psi(s)=[x_s]_{R^M}.
\]

\item \textit{Well-defined}: if $x$ and $y$ both satisfy
\[
\DBAR(s_0,x)=\DBAR(s_0,y)=s,
\]
then by definition of $R^M$ we have $x R^M y$. Hence
\[
[x]_{R^M}=[y]_{R^M},
\]
so $\psi(s)$ does not depend on which string reaching $s$ was chosen.

\item \textit{Homomorphism}: first,
\[
\psi(s_0)=[\lambda]_{R^M},
\]
because $\lambda$ reaches $s_0$. Now let $s\in S^c$ and $a\in\Sigma$. Since
\[
\DBAR(s_0,x_sa)=\delta(\DBAR(s_0,x_s),a)=\delta(s,a),
\]
the string $x_sa$ reaches $\delta(s,a)$, so
\[
\psi(\delta(s,a))=[x_sa]_{R^M}
=\delta_{R^M}([x_s]_{R^M},a)
=\delta_{R^M}(\psi(s),a).
\]
Also,
\[
s\in F
\Leftrightarrow \DBAR(s_0,x_s)\in F
\Leftrightarrow x_s\in L(M)
\Leftrightarrow [x_s]_{R^M}\in F_{R^M}
\Leftrightarrow \psi(s)\in F_{R^M}.
\]
Thus $\psi$ preserves start state, transitions, and finality.

\item \textit{Bijection}: to show injectivity, suppose $\psi(s)=\psi(t)$. Then
\[
[x_s]_{R^M}=[x_t]_{R^M},
\]
so $x_s R^M x_t$. By definition of $R^M$,
\[
\DBAR(s_0,x_s)=\DBAR(s_0,x_t),
\]
hence $s=t$. For surjectivity, let $[x]_{R^M}$ be any state of $A_{R^M}$.
Then
\[
s=\DBAR(s_0,x)\in S^c,
\]
and by definition of $\psi$,
\[
\psi(s)=[x]_{R^M}.
\]

\item $M^c\cong A_{R^M}$. $\square$
\end{enumerate}

%%------------------------------------------------------------
\item \textbf{For the machine $A$ in Figure 3.10(a), find $E_A$, $L(A)$, $A_{R_{L(A)}}$, $R_{L(A)}$, and $A/_{E_A}$.}

\begin{enumerate}
\item
\[
E_{0A}=\{\{q_0,q_1,q_3,q_6\},\{q_2,q_4,q_5\}\}.
\]
Refining once gives
\[
E_{1A}=\{\{q_0,q_1\},\{q_3,q_6\},\{q_2\},\{q_4,q_5\}\}.
\]
Refining again gives
\[
E_{2A}=\{\{q_0\},\{q_1\},\{q_2\},\{q_3,q_6\},\{q_4,q_5\}\}.
\]
This is already stable, so $E_A=E_{2A}$.

\item The language is
\[
L(A)=\{\lambda,1\}\cup 00\Sigma^* \cup 11\Sigma^*.
\]
Indeed, $\lambda$ is accepted at $q_0$, the string $1$ is accepted at $q_1$,
and every string of length at least $2$ is accepted exactly when its first two
symbols are equal.

\item $A_{R_{L(A)}}$ has five states, represented by the classes of
$\lambda,\ 1,\ 0,\ 00,$ and $01$. Writing these as
$[\lambda], [1], [0], [00], [01]$, the transitions are
\[
\begin{array}{c|cc}
 & 0 & 1\\\hline
{}[\lambda] & [0] & [1]\\
{}[1] & [01] & [00]\\
{}[0] & [00] & [01]\\
{}[00] & [00] & [00]\\
{}[01] & [01] & [01]
\end{array}
\]
with final states $[\lambda], [1], [00]$. Here $[11]=[00]$ and
$[10]=[01]$.

\item The right-congruence classes of $R_{L(A)}$ may be written as
\[
[\lambda],\qquad [1],\qquad [0],\qquad [00]=[11],\qquad [01]=[10].
\]
Equivalently, they are
\[
\{\lambda\},\quad \{1\},\quad \{0\},\quad \{x:|x|\geq 2\text{ and the first two symbols are equal}\},
\]
and
\[
\{x:|x|\geq 2\text{ and the first two symbols are different}\}.
\]

\item $A/_{E_A}$ has states
\[
[q_0],\ [q_1],\ [q_2],\ [q_3]=[q_6],\ [q_4]=[q_5],
\]
with $[q_0]$ initial and final states $[q_0], [q_1], [q_3]=[q_6]$. Its
transition table is
\[
\begin{array}{c|cc}
 & 0 & 1\\\hline
{}[q_0] & [q_2] & [q_1]\\
{}[q_1] & [q_4]=[q_5] & [q_3]=[q_6]\\
{}[q_2] & [q_3]=[q_6] & [q_4]=[q_5]\\
{}[q_3]=[q_6] & [q_3]=[q_6] & [q_3]=[q_6]\\
{}[q_4]=[q_5] & [q_4]=[q_5] & [q_4]=[q_5]
\end{array}
\]
so $A/_{E_A}\cong A_{R_{L(A)}}$.
\end{enumerate}

%%------------------------------------------------------------
\item \textbf{Exercises 3.24--3.26.}

\begin{enumerate}
\item \textbf{Machine $B$ of Figure 3.10(b).}

\begin{enumerate}
\item
\[
E_{0B}=\{\{t_0\},\{t_1,t_2,t_3\}\}.
\]
All three nonfinal states have the same $0$- and $1$-transitions into the
nonfinal class, so $E_{1B}=E_{0B}$. Hence $E_B=E_{0B}$.

\item $L(B)=\{\lambda\}$.

\item $A_{R_{L(B)}}$ is the two-state DFA with states $[\lambda]$ and
$[\mathbf{0}]$ (equivalently, $[\mathbf{1}]$), where $[\lambda]$ is initial
and final, $[\mathbf{0}]$ is nonfinal, and both letters send
$[\mathbf{0}]$ to itself.

\item $R_{L(B)}$ has exactly two classes:
\[
[\lambda]=\{\lambda\}, \qquad [\mathbf{0}]=\Sigma^+.
\]

\item $B/_{E_B}$ has states $[t_0]$ and $[t_1]=[t_2]=[t_3]$, with
$[t_0]$ initial and final, and
\[
\delta([t_0],0)=\delta([t_0],1)=[t_1]=[t_2]=[t_3],
\]
while both letters fix $[t_1]=[t_2]=[t_3]$. Thus
$B/_{E_B}\cong A_{R_{L(B)}}$.
\end{enumerate}

\item \textbf{Machine $C$ of Figure 3.10(c).}

\begin{enumerate}
\item
\[
E_{0C}=\{\{q_1,q_4\},\{q_0,q_2,q_3,q_5\}\}.
\]
Refining once gives
\[
E_{1C}=\{\{q_1\},\{q_4\},\{q_0,q_2\},\{q_3,q_5\}\}.
\]
This is already stable, so $E_C=E_{1C}$.

\item The language is
\[
L(C)=1(1\cup 01)^*,
\]
that is, the strings that begin with $1$, end with $1$, and never contain the
substring $00$.

\item $A_{R_{L(C)}}$ has three states, represented by the classes of
$\lambda,\ 1,$ and $0$. Writing them as $[\lambda], [1], [0]$, the transition
table is
\[
\begin{array}{c|cc}
 & 0 & 1\\\hline
{}[\lambda] & [0] & [1]\\
{}[1] & [\lambda] & [1]\\
{}[0] & [0] & [0]
\end{array}
\]
with only $[1]$ final. Here $[10]=[\lambda]$, while $[11]=[1]$ and
$[00]=[01]=[0]$.

\item The classes of $R_{L(C)}$ can be represented by
$[\lambda], [1], [0]$. Concretely,
\[
[1]=1(1\cup 01)^*,
\]
\[
[\lambda]=\{\lambda\}\cup 1(1\cup 01)^*0,
\]
and
\[
[0]=\Sigma^* - \big([\lambda]\cup[1]\big).
\]

\item $C/_{E_C}$ has the four states
\[
[q_1],\ [q_4],\ [q_0]=[q_2],\ [q_3]=[q_5].
\]
The initial state is $[q_0]=[q_2]$, the final states are $[q_1]$ and $[q_4]$,
and the transition table is
\[
\begin{array}{c|cc}
 & 0 & 1\\\hline
{}[q_0]=[q_2] & [q_3]=[q_5] & [q_1]\\
{}[q_1] & [q_0]=[q_2] & [q_1]\\
{}[q_3]=[q_5] & [q_3]=[q_5] & [q_3]=[q_5]\\
{}[q_4] & [q_3]=[q_5] & [q_3]=[q_5]
\end{array}
\]
The state $[q_4]$ is disconnected, so $C/_{E_C}$ is not connected; its
connected part is isomorphic to $A_{R_{L(C)}}$.
\end{enumerate}

\item \textbf{Machine $D$ of Figure 3.10(d).}

\begin{enumerate}
\item
\[
E_{0D}=\{\{S_1,S_4,S_7\},\{S_0,S_2,S_3,S_5,S_6,S_8,S_9\}\}.
\]
Refining once gives
\[
E_{1D}=\{\{S_1,S_4\},\{S_7\},\{S_0,S_3\},\{S_2,S_5,S_6\},\{S_8,S_9\}\}.
\]
Refining again gives
\[
E_{2D}=\{\{S_1,S_4\},\{S_7\},\{S_0,S_3\},\{S_2,S_5\},\{S_6\},\{S_8\},\{S_9\}\}.
\]
This is already stable, so $E_D=E_{2D}$.

\item From the initial state $S_0$, only the classes
$[S_0]=[S_3]$, $[S_1]=[S_4]$, and $[S_2]=[S_5]$ are reachable. Therefore
$L(D)$ consists exactly of those strings over $\{a,b\}$ that contain at least
one $a$ and either end in $b$ or are of the form $b^*a$. Equivalently,
\[
L(D)=b^*a(\{a,b\}^*b\cup\{\lambda\}).
\]

\item $A_{R_{L(D)}}$ has three states represented by $[\lambda]$, $[a]$, and
$[aa]$, with transition table
\[
\begin{array}{c|cc}
 & a & b\\\hline
{}[\lambda] & [a] & [\lambda]\\
{}[a] & [aa] & [a]\\
{}[aa] & [aa] & [a]
\end{array}
\]
and only $[a]$ final.

\item The classes of $R_{L(D)}$ may be described as
\[
[\lambda]=b^*,
\]
\[
[a]=L(D),
\]
and
\[
[aa]=\{x\in\{a,b\}^*: x\text{ ends with }a\text{ and contains at least two }a\text{'s}\}.
\]

\item $D/_{E_D}$ has the seven states
\[
[S_0]=[S_3],\ [S_1]=[S_4],\ [S_2]=[S_5],\ [S_6],\ [S_7],\ [S_8],\ [S_9].
\]
Its initial state is $[S_0]=[S_3]$, its final states are $[S_1]=[S_4]$ and
$[S_7]$, and its transition table is
\[
\begin{array}{c|cc}
 & a & b\\\hline
{}[S_0]=[S_3] & [S_1]=[S_4] & [S_0]=[S_3]\\
{}[S_1]=[S_4] & [S_2]=[S_5] & [S_1]=[S_4]\\
{}[S_2]=[S_5] & [S_2]=[S_5] & [S_1]=[S_4]\\
{}[S_6] & [S_6] & [S_7]\\
{}[S_7] & [S_7] & [S_7]\\
{}[S_8] & [S_0]=[S_3] & [S_9]\\
{}[S_9] & [S_2]=[S_5] & [S_2]=[S_5]
\end{array}
\]
The connected part of this quotient is the three-state machine
$A_{R_{L(D)}}$ above.
\end{enumerate}
\end{enumerate}

%%------------------------------------------------------------
\item \textbf{Prove: if $A,B$ are reduced, connected, equivalent DFAs then $A\cong B$.}

\begin{enumerate}
\item Since $A$ is connected, for each $s\in S_A$ choose a string $x_s$ with
$\DBAR_A(s_{0A},x_s)=s$. Define
\[
\psi(s)=\DBAR_B(s_{0B},x_s).
\]

\item To show this is well defined, suppose $x$ and $y$ both satisfy
$\DBAR_A(s_{0A},x)=\DBAR_A(s_{0A},y)=s$. Then for every $z\in\Sigma^*$,
\[
xz\in L(A)\Leftrightarrow yz\in L(A),
\]
because $xz$ and $yz$ continue from the same state $s$ in $A$. Since
$L(A)=L(B)$, it follows that
\[
xz\in L(B)\Leftrightarrow yz\in L(B)
\]
for every $z$. Therefore the states $\DBAR_B(s_{0B},x)$ and
$\DBAR_B(s_{0B},y)$ have the same future language in $B$. Because $B$ is
reduced, these states must be equal. Hence $\psi$ is well defined.

\item First,
\[
\psi(s_{0A})=\DBAR_B(s_{0B},\lambda)=s_{0B}.
\]
Now let $s\in S_A$ and $a\in\Sigma$. Since
\[
\DBAR_A(s_{0A},x_sa)=\delta_A(\DBAR_A(s_{0A},x_s),a)=\delta_A(s,a),
\]
we may use $x_sa$ as a string reaching $\delta_A(s,a)$. Therefore
\[
\psi(\delta_A(s,a))=\DBAR_B(s_{0B},x_sa)=\delta_B(\DBAR_B(s_{0B},x_s),a)
=\delta_B(\psi(s),a).
\]
Also,
\[
s\in F_A \Leftrightarrow x_s\in L(A) \Leftrightarrow x_s\in L(B)
\Leftrightarrow \psi(s)\in F_B.
\]
So $\psi$ is a homomorphism.

\item \textit{Injective}: suppose $\psi(s)=\psi(t)$. Then for every
$z\in\Sigma^*$,
\[
x_sz\in L(B)\Leftrightarrow x_tz\in L(B),
\]
because $x_s$ and $x_t$ reach the same state in $B$. Since $L(A)=L(B)$, we also
have
\[
x_sz\in L(A)\Leftrightarrow x_tz\in L(A)
\]
for every $z$. Thus $s$ and $t$ have the same future language in $A$. Because
$A$ is reduced, $s=t$.

\textit{Surjective}: let $u\in S_B$. Since $B$ is connected, there exists
$y\in\Sigma^*$ such that $\DBAR_B(s_{0B},y)=u$. Let
$s=\DBAR_A(s_{0A},y)$. Then $y$ is one of the strings reaching $s$, so by the
definition of $\psi$ and well-definedness,
\[
\psi(s)=\DBAR_B(s_{0B},y)=u.
\]
Hence $\psi$ is surjective.
\end{enumerate}

%%------------------------------------------------------------
\item \textbf{Show the transition from line 3 to 4 in Theorem 3.1(6) is not $\Leftrightarrow$.}

In that step, we used $\DBAR_A(s,x)\in F_A\Rightarrow\mu(\DBAR_A(s,x))\in F_B$. The converse would require $\mu(s)\in F_B\Rightarrow s\in F_A$, which is the definition of a homomorphism (both directions). If $\mu$ is only required to satisfy $s\in F_A\Rightarrow\mu(s)\in F_B$ (not $\Leftarrow$), then we only get $L(A)\subseteq L(B)$, not equality. Example: $A$ has states $\{q_0,q_1\}$, $F_A=\{q_1\}$; $B$ has $\{p_0\}$, $F_B=\{p_0\}$ (all final); $\mu(q_i)=p_0$. Then $L(A)\subset L(B)=\Sigma^*$.

%%------------------------------------------------------------
\item \textbf{Supply reasons for equivalences in proof of Theorem 3.8.}

Reading down the proof of Theorem 3.8:
\begin{enumerate}
\item
\[
sE_{i+1A}t \Leftrightarrow
(\forall x\in\Sigma^* \ni |x|\leq i+1)(\DBAR(s,x)\in F \Leftrightarrow \DBAR(t,x)\in F)
\]
is just the definition of $E_{i+1A}$.

\item Splitting this into the strings of length at most $i$ and the strings of
length exactly $i+1$ is valid because those two cases exhaust the condition
$|x|\leq i+1$.

\item Replacing ``$|y|=i+1$'' by ``$1\leq |y|\leq i+1$'' is valid because the
strings of lengths $1,\dots,i$ are already handled by the first conjunct.

\item Passing from nonempty $y$ with $|y|\leq i+1$ to strings of the form
$\mathbf{a}x$ with $\mathbf{a}\in\Sigma$ and $|x|\leq i$ uses the fact that
every nonempty string can be written uniquely in that form.

\item
\[
(\forall x\in\Sigma^* \ni |x|\leq i)(\DBAR(s,x)\in F \Leftrightarrow \DBAR(t,x)\in F)
\Leftrightarrow sE_{iA}t
\]
is the definition of $E_{iA}$.

\item
\[
\DBAR(s,\mathbf{a}x)=\DBAR(\delta(s,\mathbf{a}),x)
\]
and
\[
\DBAR(t,\mathbf{a}x)=\DBAR(\delta(t,\mathbf{a}),x)
\]
are both instances of the definition of the extended transition function
$\DBAR$.

\item Finally,
\[
(\forall x\in\Sigma^* \ni |x|\leq i)(\DBAR(\delta(s,\mathbf{a}),x)\in F
\Leftrightarrow \DBAR(\delta(t,\mathbf{a}),x)\in F)
\Leftrightarrow \delta(s,\mathbf{a})E_{iA}\delta(t,\mathbf{a})
\]
is again the definition of $E_{iA}$.
\end{enumerate}

%%------------------------------------------------------------
\item \textbf{Minimize the machine of Figure 3.3.}

Let
\[
E_{0A}=\{\{S_0,S_5\},\{S_1,S_2,S_3,S_4\}\}.
\]
Refining once gives
\[
E_{1A}=\{\{S_0,S_5\},\{S_1,S_2\},\{S_3,S_4\}\},
\]
and this is already stable. Thus the minimal machine has the three states
\[
A=[S_0]=[S_5],\qquad B=[S_1]=[S_2],\qquad C=[S_3]=[S_4],
\]
with $A$ initial and final, and transition table
\[
\begin{array}{c|cc}
 & a & b\\\hline
A & B & B\\
B & C & C\\
C & A & A
\end{array}
\]
So the machine minimizes to a $3$-state DFA in which both letters move one step
forward around the cycle $A\to B\to C\to A$. $\square$

%%------------------------------------------------------------
\item \textbf{Examples where $A$ not reduced, $A^c$ not/is reduced.}

\begin{enumerate}
\item $A$ not reduced, $A^c$ not reduced: take
\[
A=\langle\{\pmb{a}\},\{q_0,q_1\},q_0,\delta,\{q_0,q_1\}\rangle,
\]
where
\[
\delta(q_0,\pmb{a})=q_1,\qquad \delta(q_1,\pmb{a})=q_1.
\]
Both states are reachable, so $A^c=A$. Also $q_0 E_A q_1$, because from either
state every continuation is accepted. Thus $A$ is not reduced, and neither is
$A^c$.

\item $A$ not reduced, $A^c$ is reduced: take
\[
A=\langle\{\pmb{a}\},\{q_0,q_1,q_2\},q_0,\delta,\{q_1,q_2\}\rangle,
\]
where
\[
\delta(q_0,\pmb{a})=q_1,\qquad \delta(q_1,\pmb{a})=q_1,\qquad
\delta(q_2,\pmb{a})=q_2.
\]
The state $q_2$ is unreachable, and $q_1 E_A q_2$ because from either state
every continuation is accepted. So $A$ is not reduced. However,
\[
A^c=\langle\{\pmb{a}\},\{q_0,q_1\},q_0,\delta|_{S^c\times\Sigma},\{q_1\}\rangle,
\]
and $q_0$ and $q_1$ are not equivalent in $A^c$ because $\lambda$ is rejected
from $q_0$ but accepted from $q_1$. Hence $A^c$ is reduced.
\end{enumerate}

%%------------------------------------------------------------
\item \textbf{Prove $\cong$ is an equivalence relation on DFAs.}

\begin{enumerate}
\item If $g: A\to B$ and $f: B\to C$ are isomorphisms, then $f\circ g: A\to C$ satisfies all isomorphism conditions (since both are bijective homomorphisms and composition of bijections is bijective, composition of homomorphisms is a homomorphism).

\item Symmetry: if $f: A\to B$ is an isomorphism (bijective homomorphism), then $f^{-1}: B\to A$ is also an isomorphism: $f^{-1}(s_{0B})=s_{0A}$ (since $f(s_{0A})=s_{0B}$); $f^{-1}(\delta_B(t,a))=f^{-1}(\delta_B(f(s),a))=f^{-1}(f(\delta_A(s,a)))=\delta_A(s,a)=\delta_A(f^{-1}(t),a)$; $t\in F_B\Leftrightarrow f^{-1}(t)\in F_A$.

\item Reflexivity: $\mathrm{id}: A\to A$ is an isomorphism.

\item From (a),(b),(c): $\cong$ is an equivalence relation. $\square$
\end{enumerate}

%%------------------------------------------------------------
\item \textbf{Homomorphism is not an equivalence relation on DFAs.}

Homomorphism is not symmetric. Let $\Sigma=\{a\}$, and let $A$ be the
one-state DFA with state $p$ (initial, nonfinal) and $\delta_A(p,a)=p$, so
$L(A)=\emptyset$. Let $B$ have two states $q_0,q_1$, with $q_0$ initial and
nonfinal, $q_1$ final and unreachable, and both states looping on $a$.

There is a homomorphism $\mu:A\to B$ given by $\mu(p)=q_0$: initials match,
transitions are preserved, and $p\in F_A \Leftrightarrow q_0\in F_B$.

However, there is no homomorphism $\nu:B\to A$, because $q_1\in F_B$ would have
to imply $\nu(q_1)\in F_A$, and $A$ has no final state. Therefore homomorphism
is not an equivalence relation on DFAs. $\square$

%%------------------------------------------------------------
\item \textbf{For $R$ and $R_L$ from Theorem 2.2, show homomorphism from $A_R$ to $A_{R_L}$.}

$R$ refines $R_L$: each $R$-class is contained in some $R_L$-class. Define
\[
\mu:S_{A_R}\to S_{A_{R_L}}
\qquad\text{by}\qquad
\mu([x]_R)=[x]_{R_L}.
\]
This is well defined: if $[x]_R=[y]_R$, then $xRy$, and since $R$ refines
$R_L$, also $xR_Ly$, so $[x]_{R_L}=[y]_{R_L}$.

Now
\[
\mu(s_{0,R})=\mu([\lambda]_R)=[\lambda]_{R_L}=s_{0,R_L}.
\]
Also, for any $a\in\Sigma$,
\[
\mu(\delta_R([x]_R,a))
=\mu([xa]_R)
=[xa]_{R_L}
=\delta_{R_L}([x]_{R_L},a)
=\delta_{R_L}(\mu([x]_R),a).
\]
Finally,
\[
[x]_R\in F_R
\Leftrightarrow x\in L
\Leftrightarrow [x]_{R_L}\in F_{R_L}
\Leftrightarrow \mu([x]_R)\in F_{R_L}.
\]
Therefore $\mu$ is a homomorphism from $A_R$ to $A_{R_L}$. $\square$

%%------------------------------------------------------------
\item \textbf{Prove: if there is a homomorphism $A\to B$ then $R^A$ refines $R^B$.}

$\mu: A\to B$ homomorphism. If $x R^A y$: $\DBAR_A(s_0,x)=\DBAR_A(s_0,y)$. Then $\DBAR_B(s_0,x)=\mu(\DBAR_A(s_0,x))=\mu(\DBAR_A(s_0,y))=\DBAR_B(s_0,y)$, so $x R^B y$. $\square$

%%------------------------------------------------------------
\item \textbf{Prove: if $A\cong B$ then $R^A=R^B$.}

\begin{enumerate}
\item Let $\mu:A\to B$ be an isomorphism. If $xR^A y$, then
\[
\DBAR_A(s_{0_A},x)=\DBAR_A(s_{0_A},y).
\]
Applying $\mu$ to both sides gives
\[
\mu(\DBAR_A(s_{0_A},x))=\mu(\DBAR_A(s_{0_A},y)).
\]
Because $\mu$ preserves start states and transitions, the left side is exactly
$\DBAR_B(s_{0_B},x)$ and the right side is exactly $\DBAR_B(s_{0_B},y)$. Hence
$xR^B y$. Therefore $R^A\subseteq R^B$.

Now let $xR^B y$. Since $\mu^{-1}:B\to A$ is also an isomorphism, the inclusion
proved above with $A$ and $B$ interchanged gives $xR^A y$. Hence
$R^B\subseteq R^A$. Therefore
$R^A=R^B$.

\item Equivalently, for every word $x$,
\[
\mu(\DBAR_A(s_{0_A},x))=\DBAR_B(s_{0_B},x),
\]
and this can be proved by induction on $|x|$. The base case $x=\lambda$ is
\[
\mu(\DBAR_A(s_{0_A},\lambda))=\mu(s_{0_A})=s_{0_B}=\DBAR_B(s_{0_B},\lambda).
\]
For the step, write $x=\mathbf{a}w$. Then
\[
\mu(\DBAR_A(s_{0_A},\mathbf{a}w))
=\mu(\DBAR_A(\delta_A(s_{0_A},\mathbf{a}),w))
=\DBAR_B(\mu(\delta_A(s_{0_A},\mathbf{a})),w),
\]
by the induction hypothesis, and since $\mu$ preserves transitions this equals
\[
\DBAR_B(\delta_B(s_{0_B},\mathbf{a}),w)=\DBAR_B(s_{0_B},\mathbf{a}w).
\]
Because $\mu$ is injective,
\[
\DBAR_A(s_{0_A},x)=\DBAR_A(s_{0_A},y)
\Leftrightarrow
\DBAR_B(s_{0_B},x)=\DBAR_B(s_{0_B},y),
\]
so again $R^A=R^B$.
\end{enumerate}

%%------------------------------------------------------------
\item \textbf{$A$ not homomorphic to $B$, $R^A=R^B$.}

\begin{enumerate}
\item $L(A)=L(B)$: let $\Sigma=\{a\}$. Let $B$ be the minimal $2$-state DFA for
even-length strings: $q_0$ is initial and final, $q_1$ is nonfinal, and
$\delta(q_0,a)=q_1$, $\delta(q_1,a)=q_0$. Let $A$ have the same reachable part,
plus an unreachable final state $r$ with $\delta(r,a)=r$.

Then $L(A)=L(B)$ and $R^A=R^B$, because $R^A$ only depends on the states reached
from the initial state, and $A$ and $B$ have the same reachable part. However,
there is no homomorphism $A\to B$: since $r$ is final, its image would have to
be the final state $q_0$, but then
\[
\mu(\delta_A(r,a))=\mu(r)=q_0
\]
while
\[
\delta_B(\mu(r),a)=\delta_B(q_0,a)=q_1,
\]
a contradiction.

\item $L(A)\neq L(B)$: let $A$ again be the even-length machine, and let $B$ be
the same transition graph but with final state $q_1$ instead of $q_0$, so
$L(B)$ is the odd-length strings. Then $R^A=R^B$ (both relations simply
partition $\Sigma^*$ into even and odd lengths), but $A$ is not homomorphic to
$B$ because the initial state of $A$ is final while the initial state of $B$ is
not.

\item No. Suppose $A$ and $B$ are connected, $L(A)=L(B)$, and $R^A=R^B$. Define
\[
\psi(\DBAR_A(s_{0A},x))=\DBAR_B(s_{0B},x).
\]
Because $R^A=R^B$, this is well defined: if $\DBAR_A(s_{0A},x)=\DBAR_A(s_{0A},y)$,
then $xR^Ay$, hence $xR^By$, and therefore
$\DBAR_B(s_{0B},x)=\DBAR_B(s_{0B},y)$. Since both machines are connected, the
map $\psi$ is onto: given $u\in S_B$, choose $x$ with
$\DBAR_B(s_{0B},x)=u$, and then
\[
\psi(\DBAR_A(s_{0A},x))=u.
\]
It is also one-to-one: if $\psi(s)=\psi(t)$, choose strings $x_s,x_t$ with
\[
\DBAR_A(s_{0A},x_s)=s,\qquad \DBAR_A(s_{0A},x_t)=t.
\]
Then
\[
\DBAR_B(s_{0B},x_s)=\psi(s)=\psi(t)=\DBAR_B(s_{0B},x_t),
\]
so $x_sR^Bx_t$. Since $R^A=R^B$, also $x_sR^Ax_t$, hence
\[
s=\DBAR_A(s_{0A},x_s)=\DBAR_A(s_{0A},x_t)=t.
\]
For transitions,
\[
\psi(\delta_A(\DBAR_A(s_{0A},x),a))
=\psi(\DBAR_A(s_{0A},xa))
=\DBAR_B(s_{0B},xa)
=\delta_B(\DBAR_B(s_{0B},x),a)
=\delta_B(\psi(\DBAR_A(s_{0A},x)),a).
\]
Finally, for any $x\in\Sigma^*$,
\[
\DBAR_A(s_{0A},x)\in F_A
\Leftrightarrow x\in L(A)
\Leftrightarrow x\in L(B)
\Leftrightarrow \DBAR_B(s_{0B},x)\in F_B,
\]
so $\psi$ preserves finality. Thus $\psi$ is an isomorphism
$A\cong B$, so in particular there is a homomorphism from $A$ to $B$. Such
examples therefore cannot exist.

\item Yes. The example from part (a) already works if we note that both machines
are reduced: $B$ is minimal and therefore reduced, and the extra state $r$ in
$A$ has language $\Sigma^*$, which is different from the state languages of
$q_0$ and $q_1$. So $A$ and $B$ are both reduced, $L(A)=L(B)$, and
$R^A=R^B$, but $A$ is still not homomorphic to $B$.
\end{enumerate}

%%------------------------------------------------------------
\item \textbf{Disprove: if $A$ homomorphic to $B$ then $R^A=R^B$.}

Counterexample: $A$ = 2-state DFA for $\{a^{2k}\}$, $B$ = 1-state DFA for $\Sigma^*$ (all states final). There is a homomorphism $\mu: A\to B$ mapping both states to the one state of $B$. $R^A$ has 2 classes (even/odd), $R^B$ has 1 class. $R^A\neq R^B$. $\square$

%%------------------------------------------------------------
\item \textbf{Prove or give counterexamples (homomorphisms involving $A_{R^M}$).}

\begin{enumerate}
\item Homomorphism $\psi: A_{R^M}\to M$: Yes. Define
\[
\psi([x]_{R^M})=\DBAR_M(s_0,x).
\]
This is well defined: if $[x]_{R^M}=[y]_{R^M}$, then $x R^M y$, so by
definition of $R^M$,
\[
\DBAR_M(s_0,x)=\DBAR_M(s_0,y).
\]
It is a homomorphism because
\[
\psi([\lambda]_{R^M})=\DBAR_M(s_0,\lambda)=s_0,
\]
\[
\psi(\delta_{R^M}([x]_{R^M},a))
=\psi([xa]_{R^M})
=\DBAR_M(s_0,xa)
=\delta_M(\DBAR_M(s_0,x),a)
=\delta_M(\psi([x]_{R^M}),a),
\]
and
\[
[x]_{R^M}\in F_{R^M}
\Leftrightarrow x\in L(M)
\Leftrightarrow \DBAR_M(s_0,x)\in F_M
\Leftrightarrow \psi([x]_{R^M})\in F_M.
\]

\item Isomorphism $\psi: A_{R^M}\to M$: Not always. Only if $M$ is connected (since $A_{R^M}\cong M^c$, isomorphism only if $M=M^c$, i.e., $M$ connected).

\item Homomorphism $\psi: M\to A_{R^M}$: Not always. The reachable part of $M$ maps naturally to $A_{R^M}$ (send each reachable state $s$ to $[ x_s ]_{R^M}$ where $x_s$ is any string reaching $s$), but unreachable states may obstruct extension.

\textbf{Counterexample.} Let $\Sigma=\{\mathbf{a}\}$, $S=\{s_0,s_1,s_{d}\}$, $s_0$ initial, $F=\{s_1\}$, with
\[
  \delta(s_0,\mathbf{a})=s_1,\quad \delta(s_1,\mathbf{a})=s_1,\quad \delta(s_d,\mathbf{a})=s_0.
\]
State $s_d$ is unreachable. The accessible quotient $A_{R^M}$ has two states $q_0,q_1$ with
\[
  \delta'(q_0,\mathbf{a})=q_1,\quad \delta'(q_1,\mathbf{a})=q_1,
\]
$q_0$ initial, $F'=\{q_1\}$.

A homomorphism $\psi:M\to A_{R^M}$ must satisfy $\psi(s_0)=q_0$ (initial states correspond) and $\delta'(\psi(s),\mathbf{a})=\psi(\delta(s,\mathbf{a}))$ for every $s$.
From $s_d$: $\delta'(\psi(s_d),\mathbf{a})=\psi(\delta(s_d,\mathbf{a}))=\psi(s_0)=q_0$.
But $\delta'(q_0,\mathbf{a})=q_1\neq q_0$ and $\delta'(q_1,\mathbf{a})=q_1\neq q_0$, so no state in $A_{R^M}$ can serve as $\psi(s_d)$. The homomorphism does not exist. $\square$
\end{enumerate}

%%------------------------------------------------------------
\item \textbf{Prove: if $A$ minimal then $R^A=R_{L(A)}$.}

Let
\[
A=\langle\Sigma,S,s_0,\delta,F\rangle
\]
be minimal.

First, $A$ must be connected. If not, then its connected part $A^c$ accepts the
same language as $A$ but has fewer states, contradicting minimality.

Second, $A$ must be reduced. If not, then two distinct states are $E_A$-equivalent,
so the quotient $A/_{E_A}$ accepts the same language as $A$ and has fewer states,
again contradicting minimality.

Because $A$ is connected, every state is reached by some string, so the
$R^A$-classes are in one-to-one correspondence with the states of $A$. Hence
\[
rk(R^A)=|S|.
\]
Also, Exercise 2.57(c) says that $R^A$ refines $R_{L(A)}$, so
\[
rk(R^A)\geq rk(R_{L(A)}).
\]
On the other hand, Corollary 3.2 says that $A_{R_{L(A)}}$ is minimal for
$L(A)$. Since $A$ is also minimal for $L(A)$, the two minimal DFAs have the same
number of states:
\[
|S|=rk(R_{L(A)}).
\]
Therefore
\[
rk(R^A)=|S|=rk(R_{L(A)}).
\]
A refinement of finite rank cannot be proper if the two ranks are equal. Hence
$R^A=R_{L(A)}$. $\square$

%%------------------------------------------------------------
\item \textbf{Example where Exercise 3.40 fails if $A$ not minimal.}

Take $\Sigma=\{a\}$ and let
\[
A=\langle \{a\},\{q_0,q_1,q_2\},q_0,\delta,\{q_1,q_2\}\rangle,
\]
where
\[
\delta(q_0,a)=q_1,\qquad \delta(q_1,a)=q_2,\qquad \delta(q_2,a)=q_2.
\]
This DFA is connected, but it is not reduced, because $q_1 E_A q_2$: from
either state, every continuation is accepted. Hence $A$ is not minimal. From
the initial state, the machine accepts exactly
\[
L(A)=a^+.
\]
Therefore $R_{L(A)}$ has exactly two classes, represented by $\lambda$ and $a$.
On the other hand, $\lambda$, $a$, and $aa$ reach the three distinct states
$q_0$, $q_1$, and $q_2$, so $R^A$ has three classes. Hence
\[
R^A\neq R_{L(A)}.
\]

%%------------------------------------------------------------
\item \textbf{Example where Exercise 3.40 holds even if $A$ not minimal.}

Take the minimal $2$-state DFA for the language of even-length strings over
$\{a\}$, with states $q_0$ (initial and final) and $q_1$ (nonfinal), and add an
extra unreachable state $q_2$ with $\delta(q_2,a)=q_2$. The resulting machine
$A$ is not minimal, because it is not connected. However, $R^A$ is unchanged by
the presence of $q_2$, since no string from the initial state ever reaches
$q_2$. Therefore $R^A$ still has exactly the two classes ``even length'' and
``odd length,'' which are precisely the classes of $R_{L(A)}$. So in this
example $A$ is not minimal but $R^A=R_{L(A)}$. $\square$

%%------------------------------------------------------------
\item \textbf{Quotient by arbitrary relation $R$.}

\begin{enumerate}
\item $R=E_{0A}$: $A/_{E_{0A}}$ partitions into final and nonfinal. $\delta_{E_{0A}}$ maps $[s]_{E_{0A}}\times\Sigma\to S_{E_{0A}}$: need $[s]\ni s'$ with $\delta(s,a)$ and $\delta(s',a)$ in the same $E_{0A}$-class... not necessarily. So $A/_{E_{0A}}$ is not always well-defined. Example: $A$ = 3-state DFA, $F=\{q_1\}$, $\delta(q_0,a)=q_1$, $\delta(q_2,a)=q_0$. $E_{0A}$-classes: $\{q_0,q_2\}$ (nonfinal) and $\{q_1\}$ (final). $\delta_{E_{0A}}(\{q_0,q_2\},a)$: $\delta(q_0,a)=q_1\in\{q_1\}$ but $\delta(q_2,a)=q_0\in\{q_0,q_2\}$: two different classes! Not well-defined.

\item If $R$ refines $E_A$, then $A/_R$ is well defined exactly when $R$ is
stable under the transitions:
\[
sRt \Rightarrow \delta(s,a)R\delta(t,a)
\]
for every $a\in\Sigma$.

For such well-defined refinements $R$, the analogues of Theorems 3.4, 3.5, and
3.6 behave as follows.

\textit{Analogue of Theorem 3.4}: false in general. Take $R$ to be the identity
relation on the states of a DFA $A$ that is not reduced. Then $R$ certainly
refines $E_A$, and $A/_R=A$, which is not reduced.

\textit{Analogue of Theorem 3.5}: true. If $A$ is connected and $[s]_R$ is a
state of $A/_R$, choose $x$ with $\DBAR(s_0,x)=s$. Since $A/_R$ is well defined,
\[
\DBAR_R([s_0]_R,x)=[\DBAR(s_0,x)]_R=[s]_R,
\]
so every state class is reachable.

\textit{Analogue of Theorem 3.6}: true. Because $R$ refines $E_A$, every
$R$-class is contained in a single $E_A$-class, so no $R$-class mixes final and
nonfinal states; hence $F_R$ is well defined. Compute directly:
\[
x\in L(A/_R)
\Leftrightarrow \DBAR_R([s_0]_R,x)\in F_R
\Leftrightarrow [\DBAR(s_0,x)]_R\in F_R
\Leftrightarrow \DBAR(s_0,x)\in F
\Leftrightarrow x\in L(A).
\]
Thus $L(A/_R)=L(A)$.
\end{enumerate}

%%------------------------------------------------------------
\item \textbf{Homomorphisms between $M/_{E_M}$ and $A_{R_{L(M)}}$.}

\begin{enumerate}
\item Not always. Consider the DFA
\[
M=\langle\{a\},\{s_0,s_1,s_d\},s_0,\delta,\{s_1\}\rangle
\]
with
\[
\delta(s_0,a)=s_1,\qquad \delta(s_1,a)=s_1,\qquad \delta(s_d,a)=s_0.
\]
The state $s_d$ is unreachable. The three states have distinct future languages:
\[
L(M^{s_0})=a^+,\qquad L(M^{s_1})=a^*,\qquad L(M^{s_d})=a^{\ge 2},
\]
so $E_M$ is the identity relation and $M/_{E_M}=M$.

Now $L(M)=a^+$, and $A_{R_{L(M)}}$ is the $2$-state DFA with initial
nonfinal state $q_0$, final state $q_1$, and transitions
$\delta(q_0,a)=q_1$, $\delta(q_1,a)=q_1$.

If there were a homomorphism $\psi:M/_{E_M}\to A_{R_{L(M)}}$, then
$\psi(s_0)=q_0$ and $\psi(s_1)=q_1$. Since $s_d$ is nonfinal, $\psi(s_d)$ would
have to be $q_0$. But then
\[
\psi(\delta(s_d,a))=\psi(s_0)=q_0
\]
while
\[
\delta(\psi(s_d),a)=\delta(q_0,a)=q_1,
\]
a contradiction. So the statement is false.

\item Yes. Define
\[
\psi([x]_{R_{L(M)}})= [\DBAR_M(s_0,x)]_{E_M}.
\]
This is well defined: if $xR_{L(M)}y$, then for every $z\in\Sigma^*$,
\[
xz\in L(M)\Leftrightarrow yz\in L(M),
\]
so the states $\DBAR_M(s_0,x)$ and $\DBAR_M(s_0,y)$ have the same future
language and are therefore $E_M$-equivalent.

Now
\[
\psi([\lambda]_{R_{L(M)}})=[s_0]_{E_M},
\]
and for any $a\in\Sigma$,
\[
\psi(\delta_{R_{L(M)}}([x]_{R_{L(M)}},a))
=\psi([xa]_{R_{L(M)}})
=[\DBAR_M(s_0,xa)]_{E_M}
=[\delta(\DBAR_M(s_0,x),a)]_{E_M}
=\delta_{E_M}(\psi([x]_{R_{L(M)}}),a).
\]
Finally,
\[
[x]_{R_{L(M)}}\in F_{R_{L(M)}}
\Leftrightarrow x\in L(M)
\Leftrightarrow \DBAR_M(s_0,x)\in F
\Leftrightarrow [\DBAR_M(s_0,x)]_{E_M}\in F_{E_M}.
\]
So $\psi$ is a homomorphism from $A_{R_{L(M)}}$ to $M/_{E_M}$. $\square$
\end{enumerate}

%%------------------------------------------------------------
\item \textbf{Prove: $\exists m<\|S\|$ with $E_{mA}=E_{m+1,A}$.}

The sequence $E_{0A}\supseteq E_{1A}\supseteq\cdots$ is decreasing (each refinement removes some equivalences). Since $S$ is finite, the number of equivalence classes is bounded by $|S|$. Each strict refinement step increases the number of classes by at least 1. Starting from $E_{0A}$ (which has at most 2 classes: final/nonfinal), we can have at most $|S|-1$ strict refinement steps before reaching $E_A$. So $E_{|S|-1,A}=E_{|S|,A}=E_A$, meaning $m\leq|S|-1<|S|$. $\square$

%%------------------------------------------------------------
\item \textbf{Equivalence vs.\ isomorphism.}

\begin{enumerate}
\item False. Let $B$ be the minimal DFA for the even-length strings over
\[
\Sigma=\{\pmb{a}\},
\]
with states $r_0,r_1$, start state $r_0$, final set $\{r_0\}$, and transitions
\[
\delta_B(r_0,\pmb{a})=r_1,\qquad \delta_B(r_1,\pmb{a})=r_0.
\]
Now let $A$ be the DFA obtained by adjoining an unreachable state $r_2$:
\[
A=\langle\{\pmb{a}\},\{r_0,r_1,r_2\},r_0,\delta_A,\{r_0\}\rangle,
\]
where
\[
\delta_A(r_0,\pmb{a})=r_1,\qquad \delta_A(r_1,\pmb{a})=r_0,\qquad
\delta_A(r_2,\pmb{a})=r_2.
\]
Then $L(A)=L(B)$, so $A$ and $B$ are equivalent, but they are not isomorphic
because one has three states and the other has two.

\item True: if $A\cong B$ then in particular $L(A)=L(B)$ (by Exercise 3.11).
\end{enumerate}

%%------------------------------------------------------------
\item \textbf{Prove $C_i\subseteq C_{i+1}$.}

$C_0=\{s_0\}$, $C_{i+1}=C_i\cup\{\delta(s,a):s\in C_i,a\in\Sigma\}$. Clearly $C_i\subseteq C_{i+1}$ by definition. $\square$

%%------------------------------------------------------------
\item \textbf{Prove $C_i$ = states reachable from $s_0$ by strings of length $\leq i$.}

Let $R_i=\{\DBAR(s_0,x):|x|\leq i\}$.

\textit{Base}: $R_0=\{s_0\}=C_0$. \checkmark

\textit{Step}: $R_{i+1}=\{\DBAR(s_0,x):|x|\leq i+1\}=R_i\cup\{\delta(\DBAR(s_0,x),a):|x|\leq i,a\in\Sigma\}=C_i\cup\{\delta(s,a):s\in C_i,a\in\Sigma\}=C_{i+1}$. $\square$

%%------------------------------------------------------------
\item \textbf{Prove: $C_i=C_{i+1}\Rightarrow(\forall k)(C_i=C_{i+k})$.}

\textit{Base} $k=0$: trivial. $k=1$: given.

\textit{Step}: Assume $C_i=C_{i+k}$. Then $C_{i+k+1}=C_{i+k}\cup\{\delta(s,a):s\in C_{i+k}\}=C_i\cup\{\delta(s,a):s\in C_i\}=C_{i+1}=C_i$. $\square$

%%------------------------------------------------------------
\item \textbf{Prove: for a DFA $A$, $C_i=C_{i+1}\Rightarrow C_i=S^c$.}

$S^c=\bigcup_{k\geq 0}C_k$. If $C_i=C_{i+1}$, then by Exercise 3.49, $C_i=C_{i+k}$ for all $k\geq 0$. So $S^c=\bigcup_k C_k=C_i$ (since $C_k\subseteq C_i$ for $k\leq i$ by $C_i\supseteq C_{i-1}\supseteq\cdots$, and $C_k=C_i$ for $k\geq i$). $\square$

%%------------------------------------------------------------
\item \textbf{Prove: $\exists i$ with $C_i=C_{i+1}$.}

Since $S$ is finite, the sequence $C_0\subseteq C_1\subseteq\cdots$ is non-decreasing and bounded above by $S$. So it must stabilize: $\exists i$ with $C_i=C_{i+1}$. $\square$

%%------------------------------------------------------------
\item \textbf{Prove: $\exists i\leq\|S\|$ with $C_i=S^c$.}

Each strict inclusion $C_i\subsetneq C_{i+1}$ adds at least one new state. Starting from $|C_0|=1$, after at most $|S|-1$ strict inclusions we reach $C_{|S|-1}\supseteq\cdots$, which must equal $C_{|S|}$ (no room to grow beyond $|S|$ states). So $i\leq|S|-1\leq|S|$ and $C_i=S^c$. $\square$

%%------------------------------------------------------------
\item \textbf{Argue the $C_i$ algorithm is correct and terminates.}

Correctness: by Exercise 3.48, $C_i$ is exactly the reachable states within depth $i$. By Exercise 3.51, $C_i=S^c$ when stable. Termination: by Exercise 3.51 and 3.52, stability occurs by step $|S|$. $\square$

%%------------------------------------------------------------
\item \textbf{Give two DFAs with homomorphisms in both directions but no isomorphism.}

Let
\[
L = \{a^{2k} : k\geq 0\}.
\]
Let
\[
A=\langle\{a\},\{q_0,q_1\},q_0,\delta_A,\{q_0\}\rangle,
\]
where
\[
\delta_A(q_0,a)=q_1,\qquad \delta_A(q_1,a)=q_0.
\]
This is the minimal 2-state DFA for $L$.

Let
\[
B=\langle\{a\},\{p_0,p_1,p_2\},p_0,\delta_B,\{p_0,p_2\}\rangle,
\]
where
\[
\delta_B(p_0,a)=p_1,\qquad \delta_B(p_1,a)=p_0,\qquad \delta_B(p_2,a)=p_1.
\]
Here $p_2$ is unreachable, and $p_0 E_B p_2$, so $B$ is a non-minimal DFA for
$L$.

There is a homomorphism $\mu: B\to A$ collapsing $p_0,p_2\mapsto q_0$ and $p_1\mapsto q_1$: this satisfies $\delta_A(\mu(p),a)=\mu(\delta_B(p,a))$ and $p\in F_B\Leftrightarrow\mu(p)\in F_A$.

There is also a homomorphism $\nu: A\to B$ sending $q_0\mapsto p_0$ and $q_1\mapsto p_1$ (the natural inclusion into the connected part of $B$): $\nu(\delta_A(q,a))=\delta_B(\nu(q),a)$ since $\delta_A(q_0,a)=q_1$ and $\delta_B(p_0,a)=p_1=\nu(q_1)$, and $q_0\in F_A\Leftrightarrow p_0\in F_B$.

But $A\not\cong B$ since they have different numbers of states.

%%------------------------------------------------------------
\item \textbf{$R$ and $Q$ right congruences of finite rank, $R$ refines $Q$, $L$ union of $Q$-classes.}

\begin{enumerate}
\item $L$ is a union of $R$-classes: since $R$ refines $Q$, each $Q$-class is a union of $R$-classes. So $L=\bigcup(\text{some }Q\text{-classes})=\bigcup(\text{corresponding }R\text{-classes})$.

\item Homomorphism $A_R\to A_Q$: Define $\mu([x]_R)=[x]_Q$ (well-defined since $R$ refines $Q$: $[x]_R=[y]_R\Rightarrow x R y\Rightarrow x Q y\Rightarrow[x]_Q=[y]_Q$). Check: $\mu(s_{0,R})=\mu([\lambda]_R)=[\lambda]_Q=s_{0,Q}$. $\mu(\delta_R([x]_R,a))=\mu([xa]_R)=[xa]_Q=\delta_Q([x]_Q,a)=\delta_Q(\mu([x]_R),a)$. $[x]_R\in F_R\Leftrightarrow x\in L\Leftrightarrow[x]_Q\in F_Q\Leftrightarrow\mu([x]_R)\in F_Q$.

\item Not in general. Take $\Sigma=\{\pmb{a}\}$, let
\[
[\lambda]_Q=\{\lambda\}, \qquad [\pmb{a}]_Q=\{x\in\Sigma^*:|x|\geq 1\},
\]
and let $R$ refine $Q$ by splitting the nonempty strings according to parity:
\[
[\lambda]_R=\{\lambda\},\qquad
[\pmb{a}]_R=\{x:|x|\text{ is odd}\},\qquad
[\pmb{aa}]_R=\{x:|x|\text{ is even and }|x|\geq 2\}.
\]
Let $L=\{\lambda\}$, which is a union of $Q$-classes. Then $A_Q$ has two states,
$q_0=[\lambda]_Q$ (start and final) and $q_1=[\pmb{a}]_Q$, with
\[
\delta_Q(q_0,\pmb{a})=q_1,\qquad \delta_Q(q_1,\pmb{a})=q_1.
\]
The DFA $A_R$ has three states, $r_0=[\lambda]_R$ (start and final),
$r_1=[\pmb{a}]_R$, and $r_2=[\pmb{aa}]_R$, with
\[
\delta_R(r_0,\pmb{a})=r_1,\qquad
\delta_R(r_1,\pmb{a})=r_2,\qquad
\delta_R(r_2,\pmb{a})=r_1.
\]

Assume there is a homomorphism $\psi:A_Q\to A_R$. Since start states must map to
start states, $\psi(q_0)=r_0$. Then
\[
\psi(q_1)=\psi(\delta_Q(q_0,\pmb{a}))
        =\delta_R(\psi(q_0),\pmb{a})
        =\delta_R(r_0,\pmb{a})
        =r_1.
\]
But also
\[
\psi(q_1)=\psi(\delta_Q(q_1,\pmb{a}))
        =\delta_R(\psi(q_1),\pmb{a})
        =\delta_R(r_1,\pmb{a})
        =r_2,
\]
which is impossible. Therefore there need not be a homomorphism from $A_Q$ to
$A_R$.
\end{enumerate}

%%------------------------------------------------------------
\item \textbf{Prove $A_Q$ is connected.}

Every state $[x]_Q$ of $A_Q$ is reachable: $\DBAR_Q([\lambda]_Q,x)=[x]_Q$. $\square$

%%------------------------------------------------------------
\item \textbf{Prove: if $A\cong B$ and $A$ connected then $B$ connected.}

Let $\mu: A\to B$ be an isomorphism. $\mu$ is surjective (bijection), so every state of $B$ is $\mu(s)$ for some $s\in S_A$. Since $A$ is connected, $s=\DBAR_A(s_0,x)$ for some $x$. Then $\mu(s)=\mu(\DBAR_A(s_0,x))=\DBAR_B(\mu(s_0),x)=\DBAR_B(s_{0B},x)$ (by Exercise 3.10). So every state of $B$ is reachable. $\square$

%%------------------------------------------------------------
\item \textbf{Prove: if $A\cong B$ and $B$ connected then $A$ connected.}

$\mu^{-1}: B\to A$ is also an isomorphism. By Exercise 3.56, $A$ is connected. $\square$

%%------------------------------------------------------------
\item \textbf{Disprove: homomorphism $A\to B$, $A$ connected $\Rightarrow$ $B$ connected.}

Take
\[
A=\langle\{\pmb{0}\},\{s\},s,\delta_A,\{s\}\rangle
\]
with $\delta_A(s,\pmb{0})=s$. This DFA is connected. Let
\[
B=\langle\{\pmb{0}\},\{t,u\},t,\delta_B,\{t\}\rangle
\]
where $\delta_B(t,\pmb{0})=t$ and $\delta_B(u,\pmb{0})=u$. The state $u$ is
unreachable from $t$, so $B$ is not connected.

Define $\mu:A\to B$ by $\mu(s)=t$. Then $\mu$ preserves the start state, it
preserves finality because $s$ and $t$ are both final, and
\[
\mu(\delta_A(s,\pmb{0}))=\mu(s)=t=\delta_B(t,\pmb{0})
=\delta_B(\mu(s),\pmb{0}).
\]
Hence $\mu$ is a homomorphism even though $B$ is not connected.

%%------------------------------------------------------------
\item \textbf{Disprove: homomorphism $A\to B$, $B$ connected $\Rightarrow$ $A$ connected.}

Take
\[
A=\langle\{\pmb{0}\},\{s_0,s_1\},s_0,\delta_A,\{s_0,s_1\}\rangle
\]
with $\delta_A(s_0,\pmb{0})=s_0$ and $\delta_A(s_1,\pmb{0})=s_1$. The state $s_1$
is unreachable, so $A$ is not connected. Let
\[
B=\langle\{\pmb{0}\},\{t\},t,\delta_B,\{t\}\rangle
\]
with $\delta_B(t,\pmb{0})=t$; this DFA is connected.

Define $\mu:A\to B$ by
\[
\mu(s_0)=t,\qquad \mu(s_1)=t.
\]
Then $\mu$ preserves the start state, it preserves finality because all these
states are final, and
\[
\mu(\delta_A(s_i,\pmb{0}))=\mu(s_i)=t=\delta_B(t,\pmb{0})
=\delta_B(\mu(s_i),\pmb{0})
\]
for $i=0,1$. So there is a homomorphism from $A$ to connected $B$ even though
$A$ is not connected.

%%------------------------------------------------------------
\item \textbf{$A_{R^A}\cong A^c$: define $\psi$ and prove it is an isomorphism.}

\begin{enumerate}
\item $\psi([x]_{R^A})=\DBAR(s_0,x)$: sends each $R^A$-class (= set of strings reaching the same state) to that state.

\item Well-defined: $[x]_{R^A}=[y]_{R^A}\Rightarrow\DBAR(s_0,x)=\DBAR(s_0,y)$. \checkmark

\item Homomorphism: $\psi([\lambda]_{R^A})=s_0$. $\psi([xa]_{R^A})=\DBAR(s_0,xa)=\delta(\DBAR(s_0,x),a)=\delta(\psi([x]_{R^A}),a)$. $[x]\in F_{R^A}\Leftrightarrow x\in L(A)\Leftrightarrow\DBAR(s_0,x)\in F=F^c$. \checkmark

\item Isomorphism: Injective: $\psi([x])=\psi([y])\Rightarrow\DBAR(s_0,x)=\DBAR(s_0,y)\Rightarrow[x]=[y]$. Surjective (onto $S^c$): any $s\in S^c$ has $s=\DBAR(s_0,x_s)=\psi([x_s])$.
\end{enumerate}

%%------------------------------------------------------------
\item \textbf{Prove $\psi=\mu^{-1}$ when $A,B$ connected and $\psi: A\to B$, $\mu: B\to A$ isomorphisms.}

$(\mu\circ\psi)(s)=\mu(\psi(s))$. For connected $A$: $s=\DBAR_A(s_0,x)$ for some $x$. $\psi(s)=\DBAR_B(s_{0B},x)$ (by Exercise 3.10 applied to $\psi$). $\mu(\psi(s))=\mu(\DBAR_B(s_{0B},x))=\DBAR_A(s_{0A},x)=s$ (by Exercise 3.10 applied to $\mu$). So $\mu\circ\psi=\mathrm{id}_A$, i.e., $\psi=\mu^{-1}$. $\square$

%%------------------------------------------------------------
\item \textbf{Give example for which $\psi\neq\mu^{-1}$ (not connected).}

Take $A=B$ to be the DFA
\[
\langle\{\pmb{0}\},\{q_0,q_1,q_2\},q_0,\delta,\{q_0\}\rangle
\]
where
\[
\delta(q_0,\pmb{0})=q_0,\qquad \delta(q_1,\pmb{0})=q_1,\qquad
\delta(q_2,\pmb{0})=q_2.
\]
This machine is not connected because $q_1$ and $q_2$ are unreachable.

Let $\psi:A\to B$ be the identity isomorphism:
\[
\psi(q_i)=q_i \qquad (i=0,1,2).
\]
Let $\mu:B\to A$ swap the two unreachable states:
\[
\mu(q_0)=q_0,\qquad \mu(q_1)=q_2,\qquad \mu(q_2)=q_1.
\]
Because $q_1$ and $q_2$ have exactly the same final/nonfinal status and the same
loop transition on $\pmb{0}$, $\mu$ is also an isomorphism. But
\[
\mu^{-1}=\mu\neq\psi,
\]
since, for example, $\mu^{-1}(q_1)=q_2$ while $\psi(q_1)=q_1$.

%%------------------------------------------------------------
\item \textbf{Three-state DFA with $E_{0A}$ having one class. Can $E_{0A}\neq E_{1A}$?}

$E_{0A}$ has one class iff all states are equivalent under $E_{0A}$, i.e., all are final or all are nonfinal. Example: $F=S=\{q_0,q_1,q_2\}$ (all final). Then $E_{0A}=S\times S$ (one class). Can $E_{1A}\neq E_{0A}$? $E_{1A}$: $q E_{1A} r\Leftrightarrow q E_{0A} r\wedge(\forall a)\delta(q,a) E_{0A} \delta(r,a)$. Since $E_{0A}=S\times S$, all $\delta$-targets are in the same class, so $E_{1A}=E_{0A}$. So no, if $E_{0A}$ has one class then $E_{1A}=E_{0A}$.

%%------------------------------------------------------------
\item \textbf{If $A,B$ reduced and connected and $\psi: A\to B$ is a homomorphism, must $\psi$ be an isomorphism?}

Yes. Since $\psi$ is a homomorphism, it already satisfies the start-state,
final-state, and transition conditions from Definition 3.5. We only need to show
that $\psi$ is one-to-one and onto.

To prove injectivity, assume $\psi(s)=\psi(t)$. Then certainly
$\psi(s) E_B \psi(t)$, and Exercise 3.20 gives
\[
\psi(s) E_B \psi(t)\Leftrightarrow s E_A t.
\]
Since $A$ is reduced, $s E_A t$ implies $s=t$. Thus $\psi$ is one-to-one.

To prove surjectivity, let $u\in S_B$. Because $B$ is connected, there exists
$x\in\Sigma^*$ such that $\DBAR_B(s_{0B},x)=u$. Let
$s=\DBAR_A(s_{0A},x)\in S_A$. By Exercise 3.10,
\[
\psi(s)=\psi(\DBAR_A(s_{0A},x))
=\DBAR_B(\psi(s_{0A}),x)
=\DBAR_B(s_{0B},x)
=u.
\]
So every state of $B$ lies in the image of $\psi$.

Therefore $\psi$ is a bijective homomorphism, hence an isomorphism. $\square$

%%------------------------------------------------------------
\item \textbf{Prove Corollary 3.2.}

By Theorem 2.2, $A_{R_L}$ accepts $L$. The discussion immediately preceding the
corollary shows that if $M=\langle\Sigma,S,s_0,\delta,F\rangle$ is any DFA with
$L(M)=L$, then
\[
\|S\|\geq rk(R^M)\geq rk(R_L)=\|S_{R_L}\|.
\]
Thus every DFA accepting $L$ has at least as many states as $A_{R_L}$, so
$A_{R_L}$ is a minimal DFA for $L$. $\square$

%%------------------------------------------------------------
\item \textbf{Prove $S^c=\{\DBAR(s_0,x):|x|\leq\|S\|\}$.}

Let
\[
C_0=\{s_0\},\qquad
C_{i+1}=C_i\cup\{\delta(s,\mathbf{a})\mid s\in C_i,\ \mathbf{a}\in\Sigma\}.
\]
We prove by induction on $i$ that
\[
C_i=\{\DBAR(s_0,x)\mid |x|\leq i\}.
\]
For $i=0$, both sides are $\{s_0\}$. If the claim holds for $i$, then a state is
reachable by some word of length at most $i+1$ exactly when either it is already
reachable by a word of length at most $i$, or it has the form
$\delta(\DBAR(s_0,x),\mathbf{a})$ for some $|x|\leq i$ and $\mathbf{a}\in\Sigma$.
That is exactly the definition of $C_{i+1}$.

The chain
\[
C_0\subseteq C_1\subseteq C_2\subseteq\cdots\subseteq S^c
\]
can gain a new state at most $\|S\|$ times, since there are only $\|S\|$ states
altogether. Hence for some $i\leq \|S\|$ we have $C_i=C_{i+1}$, and then no new
reachable states ever appear. Therefore $C_i=S^c$ for some $i\leq \|S\|$.

Thus every connected state has the form $\DBAR(s_0,x)$ for some word with
$|x|\leq \|S\|$, so
\[
S^c\subseteq\{\DBAR(s_0,x)\mid |x|\leq\|S\|\}.
\]
The reverse inclusion is immediate, because every state reached from $s_0$ by a
word is connected. Therefore
\[
S^c=\{\DBAR(s_0,x)\mid |x|\leq\|S\|\}. \qed
\]

%%------------------------------------------------------------
\item \textbf{Prove $s E_A t\Leftrightarrow L(A^s)=L(A^t)$.}

This is the definition of $E_A$. Formally: $s E_A t$ is defined as $L(A^s)=L(A^t)$ where $A^s=\langle\Sigma,S,s,\delta,F\rangle$. The equivalence is by definition. (The substantive content is showing $E_A$ satisfies the right congruence property and the iterative characterization from Exercises 3.3--3.4.) $\square$

%%------------------------------------------------------------
\item \textbf{Properties of initial and terminal sets.}

\begin{enumerate}
\item $\{I(A,t)\mid t\in S\}$ partitions $\Sigma^*$: Each $x\in\Sigma^*$ reaches exactly one state $\DBAR(s_0,x)=t$. So $x\in I(A,t)$ for exactly one $t$. The sets are disjoint and cover $\Sigma^*$. $\square$

\item Terminal sets need not partition $\Sigma^*$. For example, let
$A=\langle\{\pmb{a}\},\{q_0,q_1\},q_0,\delta,\{q_0,q_1\}\rangle$ with
$\delta(q_0,\pmb{a})=q_0$ and $\delta(q_1,\pmb{a})=q_1$. Then
\[
T(A,q_0)=\Sigma^*=T(A,q_1).
\]
Thus the terminal sets are not disjoint, so they do not form a partition.

\item Terminal sets can partition $\Sigma^*$. Let
$A=\langle\{\pmb{a},\pmb{b}\},\{q_0,q_1\},q_0,\delta,\{q_0\}\rangle$ where both
letters toggle the states:
\[
\delta(q_0,\pmb{a})=\delta(q_0,\pmb{b})=q_1,\qquad
\delta(q_1,\pmb{a})=\delta(q_1,\pmb{b})=q_0.
\]
Then $T(A,q_0)$ is the set of even-length strings, while $T(A,q_1)$ is the set
of odd-length strings. These two sets are disjoint and their union is all of
$\Sigma^*$, so in this example the terminal sets do form a partition.
\end{enumerate}

\end{enumerate}
